\chapter{Requirements Analysis}
\label{ch:requirements_analysis}

\section{Introduction}
This chapter provides a detailed analysis of the requirements for the Seasonal Job Matching Platform Mobile Application (\texttt{job\_seeker}). It defines the system's stakeholders, functional and non-functional requirements, and operational constraints. This analysis establishes the system boundaries and baseline parameters necessary to guide the system design and implementation phases.

\section{Stakeholders and System Actors}
The software system interacts with both human and external software entities. The primary actors identified for the mobile application are:

\begin{itemize}
    \item \textbf{Applicant (Job Seeker):} The primary human stakeholder. The applicant interacts with the mobile client to manage their profile, upload resumes, select fields of interest, search/filter seasonal job listings, submit job applications, and monitor application progress.
    \item \textbf{Backend REST API (Spring Boot System):} The supporting system actor. The backend provides secure business endpoints, executes the matching recommendations algorithm, and records database transactions.
    \item \textbf{Firebase Cloud Messaging (FCM):} The external infrastructure actor. FCM handles device registration tokens and routes background push notifications to the client.
\end{itemize}

\section{Functional Requirements}
The functional requirements specify the exact actions that the mobile application must perform. These requirements are grouped into six core operational modules as outlined in Table~\ref{tab:functional_requirements}.

\begin{table}[htbp]
\centering
\caption{System Functional Requirements by Module}
\label{tab:functional_requirements}
\begin{tabular}{|l|p{10cm}|}
\hline
\textbf{Module ID} & \textbf{Functional Requirement Description} \\ \hline
\textbf{FR-AUTH-1} & The system shall allow users to register an account by providing their full name, country, phone number, email address, and a secure password. \\ \hline
\textbf{FR-AUTH-2} & The system shall validate login credentials via a secure REST API endpoint. \\ \hline
\textbf{FR-AUTH-3} & The system shall persist the user session securely across application restarts. \\ \hline
\textbf{FR-AUTH-4} & The system shall invalidate all cached credentials and user data on user logout. \\ \hline
\textbf{FR-PROF-1} & The system shall aggregate and display user profile details including contact information, applied job lists, and favorites. \\ \hline
\textbf{FR-PROF-2} & The system shall allow the user to modify their profile data and perform optimistic local updates. \\ \hline
\textbf{FR-PROF-3} & The system shall support resume management, allowing the user to create, edit, and update sections (Education, Skills, Experience) dynamically. \\ \hline
\textbf{FR-JOB-1}  & The system shall support paginated loading of job listings (default 50 items per page) to reduce startup latency. \\ \hline
\textbf{FR-JOB-2}  & The system shall support real-time job searching by title/company and filtering by location, salary, and employment type. \\ \hline
\textbf{FR-JOB-3}  & The system shall dim already viewed job listings within the current active session. \\ \hline
\textbf{FR-APP-1}  & The system shall permit users to apply for jobs with a custom description (cover letter). \\ \hline
\textbf{FR-APP-2}  & The system shall prevent duplicate job applications by disabling the apply button and verifying existing states. \\ \hline
\textbf{FR-APP-3}  & The system shall render a chronological timeline tracking application progress (Pending, Approved, Rejected). \\ \hline
\textbf{FR-NOT-1}  & The system shall register the device's FCM token upon successful authentication. \\ \hline
\textbf{FR-NOT-2}  & The system shall listen to incoming notification streams and update in-app read statuses and badge indicators dynamically. \\ \hline
\end{tabular}
\end{table}

\section{Use Case Analysis}
The system's functional workflows are modeled using a UML Use Case diagram. Figure~\ref{fig:use_case_diagram} depicts the interaction between the primary actor (Applicant), the supporting services, and the backend system.

\begin{figure}[htbp]
\centering
\begin{tikzpicture}[scale=0.9, every node/.style={transform shape}]
    % Actors
    \node[draw, minimum width=1cm, minimum height=1.5cm] (applicant) at (0, 4) {Applicant};
    \node[draw, minimum width=1cm, minimum height=1.5cm, align=center] (api) at (10, 5) {Backend\\API};
    \node[draw, minimum width=1cm, minimum height=1.5cm, align=center] (fcm) at (10, 2) {Firebase\\FCM};

    % Boundaries
    \draw (2,-1) rectangle (8,9.5);
    \node at (5, 9.1) {\textbf{Mobile Client Boundary}};

    % Use Cases
    \node[draw, ellipse, align=center] (uc1) at (5, 8) {Register / Login};
    \node[draw, ellipse, align=center] (uc2) at (5, 6.7) {Search \& Filter\\Job Listings};
    \node[draw, ellipse, align=center] (uc3) at (5, 5.2) {View Job Details};
    \node[draw, ellipse, align=center] (uc4) at (5, 3.7) {Apply for Job};
    \node[draw, ellipse, align=center] (uc5) at (5, 2.2) {Manage Resume \&\\Fields of Interest};
    \node[draw, ellipse, align=center] (uc6) at (5, 0.7) {Track Application\\Timeline};
    \node[draw, ellipse, align=center] (uc7) at (5, -0.4) {Receive Notifications};

    % Connections (Applicant)
    \draw[->] (applicant) -- (uc1);
    \draw[->] (applicant) -- (uc2);
    \draw[->] (applicant) -- (uc3);
    \draw[->] (applicant) -- (uc4);
    \draw[->] (applicant) -- (uc5);
    \draw[->] (applicant) -- (uc6);
    \draw[->] (applicant) -- (uc7);

    % Connections (Backend/FCM)
    \draw[->] (uc1) -- (api);
    \draw[->] (uc2) -- (api);
    \draw[->] (uc3) -- (api);
    \draw[->] (uc4) -- (api);
    \draw[->] (uc5) -- (api);
    \draw[->] (uc6) -- (api);
    \draw[->] (uc7) -- (fcm);
\end{tikzpicture}
\caption{UML Use Case Diagram for the Mobile Client}
\label{fig:use_case_diagram}
\end{figure}

\section{Detailed User Stories}
To align the development steps with user needs, the following detailed user stories describe core application operations along with their strict acceptance criteria.

\subsection{User Story 1: Profile Creation and Fields of Interest Tagging}
\textbf{As an} Applicant,\\
\textbf{I want to} build a digital profile and select fields of interest (e.g., Agriculture, Tourism),\\
\textbf{So that} I can present my skills and receive customized job recommendations matching my preferences.
\begin{itemize}
    \item \textbf{Acceptance Criteria 1:} The system shall display editable form fields for Name, Country, Email, and Phone number.
    \item \textbf{Acceptance Criteria 2:} The fields of interest must be managed via a dynamic tagging chip interface. Adding or removing tags must trigger an asynchronous network request.
    \item \textbf{Acceptance Criteria 3:} Any validation error (e.g., invalid phone format) must trigger visual feedback inline.
\end{itemize}

\subsection{User Story 2: Infinite Job Scrolling and Advanced Searching}
\textbf{As an} Applicant,\\
\textbf{I want to} search job lists using titles and refine options using employment categories and salary sliders,\\
\textbf{So that} I can quickly locate relevant seasonal positions without encountering UI jank or delays.
\begin{itemize}
    \item \textbf{Acceptance Criteria 1:} The job list shall load the first 200 jobs automatically via infinite scroll pagination.
    \item \textbf{Acceptance Criteria 2:} After loading 200 job items, automatic scrolling must pause, and a ``Load More'' button must appear in the footer.
    \item \textbf{Acceptance Criteria 3:} The search queries must update the active filter provider dynamically and reset the list offset to zero.
\end{itemize}

\subsection{User Story 3: Double Application Prevention}
\textbf{As an} Applicant,\\
\textbf{I want the system to} check my application state and prevent duplicate submissions,\\
\textbf{So that} I do not accidentally apply to the same seasonal job listing multiple times.
\begin{itemize}
    \item \textbf{Acceptance Criteria 1:} Tapping the ``Apply'' button must instantly trigger an optimistic UI update, appending the job ID to the local applied set and disabling the button.
    \item \textbf{Acceptance Criteria 2:} The repository must perform a verification call to the backend user applications API before committing the application request.
    \item \textbf{Acceptance Criteria 3:} If the user is confirmed to have applied, the button must remain disabled and display the text ``Applied''.
\end{itemize}

\section{Non-Functional Requirements}
Non-functional requirements describe quality attributes, system limits, and performance goals. Table~\ref{tab:non_functional_requirements} defines these requirements and their target metrics.

\begin{table}[htbp]
\centering
\caption{System Non-Functional Requirements}
\label{tab:non_functional_requirements}
\begin{tabular}{|l|p{3.5cm}|p{6.5cm}|}
\hline
\textbf{Category} & \textbf{Requirement Label} & \textbf{Target Metric / Operational Rule} \\ \hline
\textbf{Performance} & Network Latency & Time-To-Interactive (TTI) shall be $\le 200\text{ms}$ when loading a 20-item job list page under standard 4G network parameters. \\ \hline
\textbf{Performance} & Main UI Thread Budget & Deserialization and widget inflation operations must execute within $16.6\text{ms}$ per frame to preserve $60\text{ FPS}$ UI fluidity. \\ \hline
\textbf{Resource Control} & Memory Optimization & The app shall virtualize list views via lazy builders to prevent graphic widget expansion from exceeding memory heap capacities. \\ \hline
\textbf{Security} & Secure Keychain Storage & Sensitive user auth tokens must be encrypted and saved in secure hardware storage, avoiding generic shared preferences. \\ \hline
\textbf{Security} & Data Leak Protection & Logouts must clear all credentials from storage and immediately reset/invalidate all 13 in-memory data providers to protect private sessions. \\ \hline
\textbf{Reliability} & Fault Recovery & The system shall support parallel resource batch fetching with partial success thresholds and exponential backoff retry alerts. \\ \hline
\textbf{Localization} & Multi-language Support & The UI layout must switch dynamically between English and Arabic layouts, adjusting text directions (LTR/RTL) accordingly. \\ \hline
\end{tabular}
\end{table}

\section{System Constraints}
The system implementation and deployment are governed by the following operational constraints:
\begin{itemize}
    \item \textbf{Network Dependency:} The mobile client relies entirely on external network availability. When offline, cache restrictions apply and network requests fail gracefully showing custom state templates.
    \item \textbf{Platform Limitation:} The application is built using the Flutter cross-platform SDK and must maintain consistent UI rendering and haptic feedback profiles on both iOS and Android platforms.
    \item \textbf{Backend Coupling:} API contracts are coupled with the Spring Boot REST API parameters. Any signature alterations in response payloads or endpoints must be refactored via generated models.
\end{itemize}
